\title{FlashAttention-3: Fast and Accurate Attention with Asynchrony and Low-precision}

\begin{abstract}
Attention serves as a foundational component within the Transformer architecture, constituting a key limiting factor for the scalability of large language models and tasks requiring long-context processing. \textsc{FlashAttention} introduced a GPU-optimized strategy that accelerates attention mechanisms by reducing redundant memory operations. Nonetheless, prior iterations have not harnessed the full potential of cutting-edge hardware; specifically, \textsc{FlashAttention-2} utilizes merely 35\% of the H100 GPU’s computational capacity. In this work, we introduce three principal innovations to enhance attention performance on Hopper GPUs: capitalizing on the asynchronous execution capability between Tensor Cores and TMA to (1) concurrently execute computation and data transfer using warp-specialized scheduling, (2) stagger block-level matrix multiplication with softmax evaluation, and (3) employ block quantization with non-coherent computation that exploits FP8 low-precision hardware acceleration. Our proposed solution, \textsc{FlashAttention-3}, realizes a 1.5-2.0$\times$ acceleration on H100 GPUs, achieving throughput up to 740 TFLOPs/s (representing 75\% of theoretical FP16 performance), and approaching 1.2 PFLOPs/s for FP8 computations. Moreover, we show that FP8 \textsc{FlashAttention-3} delivers a 2.6$\times$ reduction in numerical error when compared to a baseline FP8 attention implementation.
\end{abstract}

\section{Introduction}
\label{sec:intro}

Within the Transformer architecture~\citep{vaswani2017attention}, the attention mechanism is the dominant factor constraining computational efficiency, as evaluating self-attention between queries and keys results in a quadratic dependency on the sequence length. Enhancing attention mechanisms to manage substantially longer contexts offers the potential for advancements across various domains, including the analysis and comprehension of extensive document collections~\citep{guo2021longt5,shaham2022scrolls,peng2023yarn}, large-scale codebases~\citep{roziere2023code, li2023starcoder}, as well as novel modalities such as high-resolution imagery~\citep{chen2022scaling}, audio~\citep{gulati2020conformer}, and video data~\citep{ho2022video}. Further, this would benefit applications requiring persistent memory, like user modeling over extended interaction histories~\citep{sun2019bert4rec} and long-horizon agent planning~\citep{yao2022react}. Consequently, there is growing enthusiasm for developing methods to accelerate attention computations with long contexts, encompassing approximate algorithms~\citep{katharopoulos2020transformers,choromanski2020rethinking, tay2020efficient}, optimized software solutions~\citep{rabe2021self,dao2022flashattention,kwon2023efficient}, and even new architectural paradigms~\citep{peng2023rwkv,sun2023retentive,gu2023mamba}.

This study extends the advances made by \citet{dao2022flashattention}, who developed exact-attention algorithms designed by taking the GPU's execution paradigms and architectural nuances into account. In their work, \textsc{FlashAttention} was proposed as an innovative method for parallelizing attention, relying on a tiling scheme that consolidates all attention-related operations into a single GPU kernel, thereby eliminating the need for slow intermediate global memory access. The subsequent work by \citet{dao2023flashattention2} introduced \textsc{FlashAttention-2}, which refactored the original approach to include parallelization over the sequence length and process key and value matrix blocks within the forward pass inner loop, enhancing both occupancy and workload distribution on the GPU.

Despite these optimizations, our analysis finds that \textsc{FlashAttention-2} still demonstrates suboptimal efficiency on cutting-edge GPUs when compared with highly optimized GEMM kernels—for example, achieving only 35\% utilization on the Hopper H100 GPU, in contrast to 80–90\% achieved by GEMM. One contributing factor is the lack of adaptation to Hopper-specific instruction sets, as the implementation relies on instructions meant for earlier architectures like Ampere when programming the Tensor Cores. Recent studies, including ThunkerKitten~\citep{spector2024thunder} and cuDNN 9~\citep{cudnn9}, highlight that employing Hopper-tailored instructions and using tile-based programming approaches can yield considerable acceleration in attention computations and streamline implementation complexity.

At a deeper level, the \textsc{FlashAttention-2} algorithm operates within a basic synchronous execution framework, and its design does not leverage either asynchrony or low-precision arithmetic. Asynchrony is introduced primarily through the architectural choice to accelerate key operations in machine learning pipelines: specialized hardware such as Tensor Cores expedite matrix multiplication, while the Tensor Memory Accelerator (TMA) handles memory operations independently of the standard CUDA cores, which are responsible for logic, integer, and floating-point calculations. Employing low precision—like FP8 in Hopper and FP4 in Blackwell, following on from earlier adoption of FP16 (since Pascal, 2017) and BF16 (Ampere, 2020)—has been well established as a method for doubling or even quadrupling computational throughput within equivalent power and area constraints. Further discussion of Hopper's relevant hardware capabilities is presented in \cref{subsec:hardware}.
A central technical difficulty lies in adapting \textsc{FlashAttention-2} to capitalize on these architectural advancements: exploiting asynchrony involves orchestrating matmul and softmax computations such that their execution overlaps, despite inherent data dependencies; and utilizing lower-precision formats mandates rigorous handling to suppress quantization artifacts, with particular attention required for outlier activations commonly observed in LLMs~\citep{dettmers2208llm, sun2024massive}.

In pursuit of this goal, we introduce \textsc{FlashAttention-3}, which integrates and advances three novel strategies aimed at enhancing efficiency on contemporary GPU platforms.\footnote{While our findings are illustrated using NVIDIA's Hopper architecture, the proposed algorithm is applicable to any GPU system that supports strong asynchronous execution and low-precision operations.}

\begin{enumerate}

\item \textbf{Producer-Consumer asynchrony:} We introduce a warp-specialized approach to software pipelining, capitalizing on the decoupled nature of data movement operations and Tensor Core computations. This is achieved by assigning the roles of data producers and data consumers to distinct warps, thus enhancing the algorithm’s capacity to obscure both memory access and instruction issuance delays.
\item \textbf{Hiding softmax under asynchronous block-wise GEMMs:} We interleave low-throughput softmax-related non-GEMM operations—such as floating point multiply-adds and exponentials—with asynchronously launched WGMMA GEMM operations. To facilitate this, we restructure the \textsc{FlashAttention-2} procedure in order to break explicit sequential constraints between the softmax and the associated GEMMs. For instance, in our algorithm’s two-stage configuration, the softmax operates over one score matrix block while, in parallel, WGMMA asynchronously computes subsequent blocks.
\item \textbf{Hardware-accelerated low-precision GEMM:} We modify the forward pass to support execution on FP8 Tensor Cores for GEMM operations, resulting in nearly a twofold increase in observed TFLOPs/s. Accomplishing this demands adaptation to the divergent memory layout conventions expected by WGMMA when handling FP32 accumulator blocks versus FP8 operand blocks. We apply block quantization and incoherent processing strategies to alleviate any accuracy degradation attributable to the use of reduced-precision FP8 calculations.

To empirically assess the effectiveness of our approach, we conduct benchmarks of \textsc{FlashAttention-3} on the H100 SXM5 GPU across various configuration settings. Our findings indicate the following: (1) FP16 mode delivers a forward-pass speedup of 1.5-2.0$\times$ compared to \textsc{FlashAttention-2}, reaching peak throughput of up to 740 TFLOPs/s, and achieves a 1.5-1.75$\times$ increase in speed during the backward pass; (2) FP8 mode attains performance near 1.2 PFLOPs/s; and (3) for longer sequence lengths, FP16 mode consistently surpasses competing methods, while FP8 remains comparable\footnote{Specifically, FP8 mode in \textsc{FlashAttention-3} leads for head dimension 64. For head dimensions of 128 and 256, it matches the alternatives in settings without causal masking and lags slightly when causal masking is applied.} to the state-of-the-art attention implementation available in NVIDIA’s cuDNN library. In terms of numerical precision, we confirm that FP16 \textsc{FlashAttention-3} yields error rates equivalent to those of \textsc{FlashAttention-2}, and further, surpasses the accuracy of the baseline attention method due to the retention of intermediate calculations (such as softmax rescaling) in FP32 format. Additionally, when employing block quantization and incoherent processing, FP8 \textsc{FlashAttention-3} demonstrates a 2.6$\times$ enhancement in accuracy relative to conventional attention utilizing per-tensor quantization, especially in scenarios characterized by outlier features.

We have released \textsc{FlashAttention-3} as open-source software under a permissive license\footnote{\textsc{FlashAttention-3}
  is available at \texttt{https://github.com/Dao-AILab/flash-attention}}. Our aim is to further expand its accessibility by integrating it into the PyTorch and Hugging Face ecosystems, thereby providing broad utility for the research and developer communities.

\section{Background: Multi-Head Attention and GPU Characteristics}
\label{sec:background}

\subsection{Multi-Head Attention}
\label{subsec:multi_head_attn}

Let $\mathbf{Q}, \mathbf{K}, \mathbf{V} \in \mathbb{R}^{N \times d}$ be the query, key and value input sequences associated to a single head, where $N$ is the sequence length and $d$ is the head dimension. Then the attention output $\mathbf{O}$ is computed as:

\begin{equation*}
  \mathbf{S} = \alpha \mathbf{Q} \mathbf{K}^\top \in \mathbb{R}^{N \times N}, \quad \mathbf{P} = \mathrm{softmax}(\mathbf{S}) \in \mathbb{R}^{N \times N}, \quad \mathbf{O} = \mathbf{P}\mathbf{V} \in \mathbb{R}^{N \times d},
\end{equation*}

Here, the $\mathrm{softmax}$ function is computed across each row, and it is common to choose the scaling parameter as $\alpha = 1/\sqrt{d}$. To mitigate potential numerical issues arising from the exponentiation, one usually offsets $\mathbf{S}$ by $\mathrm{rowmax}(\mathbf{S})$. In the case of multi-head attention (MHA), each attention head maintains independent projections for the queries, keys, and values; this operation is parallelized over both the head and batch dimensions, generating the final output tensor.

Now let $\phi$ be a scalar loss function and let $\mathbf{d}(-) = \partial \phi / \partial (-)$ be notation for the gradient. Given the output gradient $\mathbf{dO} \in \mathbb{R}^{N \times d}$, we compute $\mathbf{dQ}$, $\mathbf{dK}$, and $\mathbf{dV}$ according to the chain rule as follows:

\begin{align*}
  \mathbf{dV} &= \mathbf{P}^\top \mathbf{dO} \in \mathbb{R}^{N \times d} \\
  \mathbf{dP} &= \mathbf{dO} \mathbf{V}^\top \in \mathbb{R}^{N \times N} \\
  \mathbf{dS} &= \mathrm{dsoftmax} (\mathbf{dP}) \in \mathbb{R}^{N \times N} \\
  \mathbf{dQ} &= \alpha \mathbf{dS} \mathbf{K} \in \mathbb{R}^{N \times d} \\
  \mathbf{dK} &= \alpha \mathbf{dS}^\top \mathbf{Q} \in \mathbb{R}^{N \times d},
\end{align*}

In this context, the relationship $\mathbf{d}s = (\mathrm{diag}(p) - p p^\top)\mathbf{d}p$ holds when $p$ is expressed as $\mathrm{softmax}(s)$, where $s$ denotes a vector. The notation $\mathrm{dsoftmax}(\mathbf{dP})$ is utilized to indicate this calculation performed on each row individually. Furthermore, this operation can be efficiently executed in parallel for all heads and batches during the backward propagation step of MHA.

\subsection{GPU hardware characteristics and execution model}
\label{subsec:hardware}

We outline the features of the GPU execution model pertinent to \textsc{FlashAttention-3}, concentrating on the NVIDIA Hopper architecture as a representative example of this model.

\paragraph{Memory hierarchy:} On GPUs, memory is structured in a hierarchy of data locations, where higher bandwidth accompanies lower capacity (\cref{tab:gpu-hierarchy})\footnote{\citet{luo2024benchmarking} lists a shared memory bandwidth of 128 bytes per clock cycle per SM; this is then scaled by 132 SMs and a boost clock rate of 1830 MHz.}. At the base of the hierarchy sits global memory (GMEM), or HBM, functioning as off-chip DRAM and providing access to all streaming multiprocessors (SMs). When data is required, it is automatically cached from GMEM into an on-chip L2 cache. Beyond this level, each SM is equipped with a limited-capacity, highly banked, programmer-controlled cache known as shared memory (SMEM). The highest-speed storage is the register file, found inside each SM.

\paragraph{Thread hierarchy:} In the GPU programming model, execution units known as threads are arranged into several logical groups. Starting with individual threads at the most granular level, threads are grouped into warps (sets of 32 threads), which are then organized into warpgroups (containing 4 consecutive warps). Multiple warpgroups form threadblocks (also referred to as cooperative thread arrays or CTAs). On some architectures, such as Hopper, threadblock clusters consist of multiple threadblocks. At the highest level, all these structures collectively comprise a grid.

The two hierarchies exhibit a strong interdependence. Threads belonging to the same CTA are executed together on a single SM, and CTAs within one cluster are collectively scheduled on the same GPC. Within a CTA, all threads can directly access SMEM, while each thread is restricted to a maximum of 256 private registers (RMEM) exclusive to itself.

\paragraph{Asynchrony and warp-specialization:}

GPUs are designed as throughput-oriented processors that utilize parallelism and asynchronous operations to mask delays associated with memory accesses and computation. In the Hopper architecture, asynchronous transfers between GMEM and SMEM are facilitated by the Tensor Memory Accelerator (TMA), a specialized piece of hardware \cite[\S7.29]{cuda}. Additionally, departing from earlier designs like Ampere, Hopper’s Tensor Core supports asynchronous operations and, via the warpgroup-wide WGMMA instruction \cite[\S9.7.14]{ptx}, can fetch its input data directly from shared memory.

Enabling asynchrony through hardware features facilitates the design of warp-specialized kernels, whereby individual warps within a CTA are assigned distinct functions, operating either as producers focused solely on data movement or as consumers dedicated exclusively to computation. This structural separation enhances the compiler's potential to construct more efficient instruction schedules \citep{warp-specialization-2011}. Furthermore, Hopper provides mechanisms for the flexible redistribution of registers among warpgroups via \verb|setmaxnreg| \cite[\S9.7.17.1]{ptx}, thereby allowing warps responsible for MMAs to access an increased portion of RMEM compared to those limited to issuing TMA instructions, which may require only a single thread.

\paragraph{Low-precision number formats:}
\label{sec:low-precision-gpu}
Contemporary GPUs are equipped with dedicated hardware components optimized for low-precision arithmetic operations. As an illustration, the WGMMA instruction is designed to leverage the FP8 Tensor Cores available on Hopper architectures, enabling twice the throughput per SM relative to operations performed in FP16 or BF16 formats.

Nonetheless, properly utilizing FP8 WGMMA requires an appreciation of the operand layout limitations it imposes. Consider performing a GEMM operation where one computes $A \times B^{\top}$, with $A$ being an $M \times K$ matrix and $B$ an $N \times K$ matrix. The $A$ or $B$ operand is described as \emph{mn-major} if its memory layout is contiguous in the outer $M$ or $N$ dimension, while it is termed \emph{k-major} if contiguity lies along the inner $K$ dimension. With FP16 WGMMA, the hardware accommodates both mn-major and k-major memory layouts for input operands stored in SMEM. In contrast, FP8 WGMMA restricts operand layouts to only the k-major format. This restriction becomes particularly problematic in cases such as attention mechanisms, where it is desirable to chain multiple GEMMs within a single kernel; conflicting data layouts between FP32 accumulators and FP8 operands hinder the ability to directly launch successive, dependent FP8 WGMMAs.

With regard to attention mechanisms, these layout constraints necessitate specific adjustments to how an FP8 algorithm is implemented, as outlined in \cref{sec:algofp8}.

\subsection{Standard Attention and Flash Attention} As outlined by \citet{dao2022flashattention}, \textbf{standard attention} refers to an attention computation on GPUs where the intermediate matrices $\mathbf{S}$ and $\mathbf{P}$ are explicitly stored in high-bandwidth memory (HBM). In contrast, the principal innovation of \textsc{FlashAttention} is its use of a localized softmax reduction, which eliminates the costly process of reading and writing these intermediate results. This technique integrates the entire attention process within a unified computational kernel. The local softmax procedure is implemented in lines \ref{code-ws:softmax_start}-\ref{code-ws:softmax_end} of the consumer mainloop in \cref{alg:flash3_wgmma_ws_only}, in combination with the normalization applied to blocks of $\mathbf{O}$. A straightforward proof that this method correctly calculates $\mathbf{O}$ is presented in \cite[\S 2.3.1]{dao2023flashattention2}.

\section{FlashAttention-3: Algorithm}
\label{sec:algo}

In this section, we present the \textsc{FlashAttention-3} algorithm. To streamline the discussion, we restrict our attention to the forward pass, while the approach for the backward pass is detailed in~\cref{sec:algo_ws_bwd}. We begin by outlining the procedure for incorporating warp-specialization along with a circular SMEM buffer into the foundational algorithm used by \textsc{FlashAttention-2}. Next, we detail the use of WGMMA asynchrony to construct an overlapped, two-stage GEMM-softmax pipeline. Lastly, we discuss the necessary adaptations for FP8, including changes for data layout compatibility and enhancement of precision through block quantization and handling of incoherent operations.

\subsection{Producer-Consumer asynchrony through warp-specialization and pingpong scheduling}
\label{sec:algo_ws}

\paragraph{Warp-specialization}
Similar to the approach in \textsc{FlashAttention-2}, the forward computation in \textsc{FlashAttention-3} can be fully parallelized across the batch dimension, head count, and the query sequence positions. Therefore, a core-thread array (CTA)-centric perspective suffices: the main computational routine processes a specific tile $\mathbf{Q}_i$ from the query matrix and produces the associated output tile $\mathbf{O}_i$. For clarity, we introduce the warp-specialization strategy by describing the use of a circular SMEM buffer in a scenario where GEMM-softmax overlap is not present. Denoting $d$ as the head size and $N$ as the full sequence length, we define the query block size $B_r$ such that $\mathbf{Q}$ is partitioned into $T_r = \lceil \frac{N}{B_r} \rceil$ submatrices, $\mathbf{Q}_1, .., \mathbf{Q}_{T_r}$.

\begin{algorithm}[H]
    \caption{\small\label{alg:flash3_wgmma_ws_only}\textsc{FlashAttention-3} forward computation \textbf{excluding} intra-consumer overlap -- CTA perspective}
    \begin{algorithmic}[1]
\REQUIRE Matrices $\mathbf{Q}_i \in \mathbb{R}^{B_r \times d}$ and $\mathbf{K}, \mathbf{V} \in \mathbb{R}^{N \times d}$ stored in HBM, key block size $B_c$ with $T_c = \lceil \frac{N}{B_c} \rceil$.
\STATE Set up pipeline object for barrier-based synchronization using an $s$-stage circular SMEM buffering scheme.
\IF {executing in the producer warpgroup}
\STATE Release a set number of preassigned registers.
\STATE Perform transfer of $\mathbf{Q}_i$ from HBM to shared memory.
\STATE On transfer completion, signal consumers that $\mathbf{Q}_i$ has been loaded.
\FOR{$0 \le j < T_c$}
    \STATE Await consumption of the $(j\,\%\,s)$ buffer stage by consumer.
    \STATE Issue transfer of $\mathbf{K}_j, \mathbf{V}_j$ from HBM into the shared memory corresponding to the $(j\,\%\,s)$ buffer stage.
    \STATE After transfer is done, notify consumers of successful load of $\mathbf{K}_j, \mathbf{V}_j$.
\ENDFOR
\ELSE \STATE Reclaim registers proportional to the number of consumer warps in use.
\STATE Initialize in on-chip memory: $\mathbf{O}_i = (0) \in \mathbb{R}^{B_r \times d}$, and for normalization, $\ell_i, m_i = (0), (-\infty) \in \mathbb{R}^{B_r}$.
\STATE Await availability of $\mathbf{Q}_i$ in shared memory.
\FOR{$0 \le j < T_c$}
\STATE Wait until $\mathbf{K}_j$ is present in shared memory.
\STATE Compute the block matrix product $\mathbf{S}_i^{(j)} = \mathbf{Q}_i \mathbf{K}_j^T$ via SS-GEMM. Commit result and block until complete. \label{alg:ws_only_gemm1}
\STATE Store previous value $m_i^{\mathrm{old}} = m_i$ and set $m_i = \mathrm{max}(m_i^{\mathrm{old}}, \mathrm{rowmax}(\mathbf{S}_i^{(j)}))$. \label{code-ws:softmax_start}
\STATE Compute unnormalized softmax numerators $\widetilde{\mathbf{P}}_i^{(j)} = \mathrm{exp}(\mathbf{S}_i^{(j)} - m_i)$ and update normalization: $\ell_i = \mathrm{exp}(m_i^{\mathrm{old}} - m_i) \ell_i + \mathrm{rowsum}(\widetilde{\mathbf{P}}_i^{(j)})$. \label{code-ws:softmax_end}
\STATE Await loading of $\mathbf{V}_j$ into shared memory.
\STATE Update output: $\mathbf{O}_i = \mathrm{diag}(\mathrm{exp}(m_i^{\mathrm{old}} - m_i))^{-1} \mathbf{O}_i + \widetilde{\mathbf{P}}_i^{(j)} \mathbf{V}_j$ via RS-GEMM. Commit and synchronize. \label{alg:ws_only_gemm2}
\STATE Mark $(j\,\%\,s)$ buffer stage as released, enabling producer operations.
\ENDFOR
\STATE Finalize output normalization: $\mathbf{O}_i = \mathrm{diag}(\ell_i)^{-1} \mathbf{O}_i$, and calculate $L_i = m_i + \log(\ell_i)$.
\STATE Persist results $\mathbf{O}_i$ and $L_i$ to HBM as block $i$ entries for $\mathbf{O}$ and $L$.
\ENDIF
\end{algorithmic}
\end{algorithm}

In our implementation of \cref{alg:flash3_wgmma_ws_only} tailored for Hopper, we utilize \verb|setmaxnreg| for managing register (de)allocations, employ TMA to load the tensors $\mathbf{Q}_i$ and $\{ \mathbf{K}_j, \mathbf{V}_j \}_{0 \leq j < T_c}$, and apply WGMMA to perform the GEMMs within the consumer mainloop. The use of the SS or RS prefix designates whether the initial operand is loaded from shared memory or the register file, respectively. When analyzing the control flow of \cref{alg:flash3_wgmma_ws_only}, it is important to observe that TMA load requests are asynchronous, allowing them to be issued independently without waiting for other loads to complete. In addition, within the producer mainloop, waits are omitted during the first $s$ iterations as the buffer is still being primed.

\paragraph{Pingpong scheduling} The asynchronous execution enabled by WGMMA, TMA, and warp-specialized strategies provides an avenue for concurrent execution, such as running softmax computations in one warpgroup while another warpgroup handles GEMM operations. This overlapping is particularly relevant given the discrepancy in throughput between matrix multiplication and other operations on present-day hardware accelerators. For instance, the H100 SXM5 GPU achieves a peak performance of 989 TFLOPS for FP16 matrix multiplications, in stark contrast to merely 3.9 TFLOPS for executing special functions like the exponential\footnote{According to the CUDA programming guide, each streaming multiprocessor (SM) can execute 16 special function operations per clock. Scaling by 132 SMs and a 1830 MHz clock rate yields an aggregate throughput of 3.9 TFLOPS for special functions.}. During an FP16 attention forward pass with a head dimension of 128, the number of floating-point operations for matrix multiplications exceeds those for exponentials by a factor of 512, yet exponentials operate at only 1/256 the throughput. Consequently, exponential operations can occupy as much as half of the computational cycles used by matrix multiplications. This imbalance becomes more pronounced in FP8 precision, where matmul throughput is doubled, but the rate of exponential computations remains unchanged.

Because exponentiation is handled by a distinct hardware unit (the multi-function unit), optimal performance would schedule the exponential computation concurrently with the matmul operations executed by the Tensor Cores. To achieve this coordination, we employ synchronization barriers (\texttt{bar.sync} instructions) so that GEMM operations (specifically, GEMM1 for $\mathbf{P} \mathbf{V}$ in one iteration and GEMM0 for $\mathbf{Q} \mathbf{K}^\top$ in the subsequent iteration) within warpgroup 1 are executed prior to those of warpgroup 2. Consequently, this arrangement allows the softmax operation of warpgroup 1 to run in parallel with the GEMM computations of warpgroup 2. Subsequently, the responsibilities alternate: while warpgroup 2 proceeds with its softmax computation, warpgroup 1 undertakes its GEMMs, resulting in a so-called ``pingpong'' execution strategy. An illustrative depiction of this method can be found in~\cref{fig:pingpong_scheduling}.

Although real-world execution deviates somewhat from the schematic representation, this form of pingpong scheduling typically yields notable performance improvements (for instance, forward pass throughput in FP16 mode with head dimension 128 and sequence length 8192 increases from approximately 570 TFLOPS to 620--640 TFLOPS).

\paragraph{Attention variants} In the case of multi-query attention~\citep{shazeer2019fast} and grouped query attention~\citep{ainslie2023gqa}, our method aligns with that of \textsc{FlashAttention-2}, modifying tensor indexing so as to prevent repeated storage of $\mathbf{K}$ and $\mathbf{V}$ within HBM.

\subsection{Intra-warpgroup overlapping GEMMs and softmax}

It is possible to interleave certain instructions from the softmax operation with those from the GEMMs, even within a single warpgroup. Here, we outline a method for achieving this overlap.

In the attention algorithm, operations within the inner loop (main loop) have sequential dependencies that impede parallelization within a single iteration. For example, (local) softmax (lines \ref{code-ws:softmax_start} to \ref{code-ws:softmax_end}) relies on the output $\mathbf{S}_i^{(j)}$ of the first GEMM, while the second GEMM takes its result $\widetilde{\mathbf{P}}_i^{(j)}$ as an operand. Indeed, the wait statements in lines \ref{alg:ws_only_gemm1} and \ref{alg:ws_only_gemm2} of \cref{alg:flash3_wgmma_ws_only} serialize the execution of softmax and GEMMs. However, we can break these dependencies by pipelining across iterations through additional buffers in registers. Pursuing this idea, we propose the following two-stage\footnote{Note that the number of stages of the overlapping scheme is bounded by, but need not equal, the number $s$ of stages in the circular SMEM buffer.} GEMM-softmax pipelining algorithm:

\begin{algorithm}[H]
    \caption{\small\label{alg:flash3_wgmma}\textsc{FlashAttention-3} Consumer Warpgroup Forward Pass}
    \begin{algorithmic}[1]
      \REQUIRE  Matrices $\mathbf{Q}_i \in \mathbb{R}^{B_r \times d}$ as well as $\mathbf{K}, \mathbf{V} \in \mathbb{R}^{N \times d}$ stored in HBM; key block size $B_c$ yielding $T_c = \lceil \frac{N}{B_c} \rceil$ blocks.
      \STATE Allocate a prescribed number of registers according to the count of consumer warps in use.
      \STATE Initialize in on-chip memory: set $\mathbf{O}_i = (0) \in \mathbb{R}^{B_r \times d}$ and vectors $\ell_i, m_i = (0), (-\infty) \in \mathbb{R}^{B_r}$.
      \STATE Await completion of loading $\mathbf{Q}_i$ and the initial key block $\mathbf{K}_0$ into shared memory.
      \STATE Execute multiplication $\mathbf{S}_{\mathrm{cur}} = \mathbf{Q}_i \mathbf{K}_0^T$ via WGMMA. Ensure commit, then synchronize.
      \STATE Release buffer stage $0$ for $\mathbf{K}$.
      \STATE Determine $m_{i}$, $\tilde{\mathbf{P}}_{\mathrm{cur}}$, and $\ell_{i}$ from $\mathbf{S}_{\mathrm{cur}}$, applying necessary rescaling to $\mathbf{O}_i$.
      \FOR{$1 \le j < T_c - 1$}
        \STATE Wait for the $j$-th key block $\mathbf{K}_j$ to finish loading into shared memory.
        \label{code:mainloop_start}
        \STATE Compute $\mathbf{S}_{\mathrm{next}} = \mathbf{Q}_i \mathbf{K}_{j}^T$ using WGMMA. Commit the operation, but do not synchronize yet. \label{code:first_wgmma}
        \STATE Wait until the value block $\mathbf{V}_{j-1}$ is loaded into shared memory.
        \STATE Update $\mathbf{O}_{i} = \mathbf{O}_{i} + \tilde{\mathbf{P}}_{\mathrm{cur}} \mathbf{V}_{j-1}$ through WGMMA, committing without synchronizing. \label{code:second_wgmma}
        \STATE Wait for completion of WGMMA product $\mathbf{Q}_i \mathbf{K}_{j}^T$.
        \STATE Using $\mathbf{S}_{\mathrm{next}}$, compute new $m_{i}$, next normalized probability $\tilde{\mathbf{P}}_{\mathrm{next}}$, and $\ell_{i}$. \label{code:softmax}
        \STATE Wait until WGMMA for $\tilde{\mathbf{P}}_{\mathrm{cur}} \mathbf{V}_{j-1}$ is finished, then apply corresponding rescaling to $\mathbf{O}_i$.
        \STATE Free up the $(j\,\%\,s)$-th buffer stage for $\mathbf{K}$ and the $(j-1\,\%\,s)$-th for $\mathbf{V}$.
        \STATE Assign $\mathbf{S}_{\mathrm{next}}$ to $\mathbf{S}_{\mathrm{cur}}$ for the following iteration.
        \label{code:mainloop_end}
      \ENDFOR \STATE Wait until the value block $\mathbf{V}_{T_c - 1}$ becomes available in shared memory.
      \STATE Perform final output update: compute
      $\mathbf{O}_{i} = \mathbf{O}_{i} + \tilde{\mathbf{P}}_{\mathrm{last}} \mathbf{V}_{T_c - 1}$ using WGMMA. Commit and synchronize before exiting.

\STATE Final step: Adjust $\mathbf{O}_{i}$ according to $m_{i}$. Determine $L_{i}$ using $m_{i}$ together with $\ell_{i}$.
Save $\mathbf{O}_{i}$ and $L_{i}$ in HBM, placing them as the $i$-th segments of $\mathbf{O}$ and $L$ respectively.
\end{algorithmic}
\end{algorithm}

\cref{alg:flash3_wgmma} serves to substitute the consumer pathway of \cref{alg:flash3_wgmma_ws_only}, thereby completing the \textsc{FlashAttention-3} workflow for FP16 computations. More generally, WGMMA is utilized to signify asynchronous GEMM throughout the discussion. During the core loop (spanning lines \ref{code:mainloop_start} to \ref{code:mainloop_end}), the execution of the second WGMMA instruction in the $j^{\text{th}}$ iteration (see line \ref{code:second_wgmma}) is scheduled concurrently with the softmax calculations from the subsequent $(j+1)^{\text{th}}$ iteration (line \ref{code:softmax}).

Although the pipelined approach depicted above promises significant theoretical speedup, several real-world considerations arise:

\paragraph{Compiler reordering} Although the pseudocode conveys an ideal sequence of instructions, in practice the compiler (NVCC) frequently modifies the ordering of instructions to optimize performance. This may interfere with the precise alternation between WGMMA and non-WGMMA operations, potentially resulting in suboptimal or unintended pipeline execution. Examination of the SASS output reveals that the compiler indeed produces overlapping code in line with expectations (see Section~\ref{sec:2-stage-sass}).

\paragraph{Register pressure} Achieving peak throughput requires minimizing register spilling, but the 2-stage pipeline introduces further register requirements for maintaining intermediate results and preserving state across pipeline phases. In particular, the design must allocate space for an additional $\mathbf{S}_{\mathrm{next}}$ within registers, increasing register usage by $B_r \times B_c \times \text{sizeof}(\text{float})$ per threadblock. Such growth in register needs may constrain the use of larger block dimensions—another widely used performance tactic—since both approaches intensify register consumption. Empirical profiling is therefore necessary to navigate these competing demands and determine the most beneficial trade-off.

\paragraph{3-stage pipelining} Building upon the 2-stage strategy, we also introduce a 3-stage version which permits overlap of the second WGMMA phase with softmax computation. This adaptation has the potential to further heighten utilization of Tensor Cores, but at the expense of additional registers to accommodate the new pipeline stage. The interplay between tile sizing and pipeline depth thus becomes even more intricate. Details of the 3-stage algorithm and experimental assessment are provided in ~\cref{sec:3-stage}.

\subsection{Low-precision with FP8}
\label{sec:algofp8}

\textbf{Efficiency: layout transformations.} Executing the forward pass of \textsc{FlashAttention-3} using FP8 precision introduces further complexities regarding layout compatibility, which are not present when utilizing FP16.

To begin, we observe that the input tensors $\mathbf{Q}$, $\mathbf{K}$, and $\mathbf{V}$ are generally organized so that the head dimension is contiguous. However, to adhere to the k-major requirement of FP8 WGMMA for the second GEMM operation, we must ensure that $\mathbf{V}$—specifically the portions of $\mathbf{V}$ brought into SMEM—is stored contiguously along the sequence length axis. As the TMA load operation does not allow us to alter the contiguous dimension, there are two viable alternatives: (1) carry out a transposition of $\mathbf{V}$ in GMEM beforehand, or (2) perform an in-kernel tile-wise transpose on $\mathbf{V}$ once it has been transferred into SMEM. Option (1) can be approached in two ways: (1a) by integrating the transpose with the epilogue of an earlier stage, such as rotary embedding, or (1b) by invoking a dedicated pre-processing transpose kernel\footnote{A highly tuned transpose kernel can closely match the device's memory bandwidth~\citep{colfax_cutlass_transpose_2024}.} that switches the strides between the sequence length and head dimensions. Nevertheless, strategy (1a) poses challenges for adoption in standard libraries, and approach (1b) incurs excessive overhead in scenarios dominated by memory throughput, such as inference workloads.

For FP8 \textsc{FlashAttention-3}, we adopt option (2) as our preferred approach. The in-kernel transpose leverages LDSM (\verb|ldmatrix|) and STSM (\verb|stmatrix|) instructions, which enable a warp of threads to collaboratively transfer data between SMEM and RMEM, operating at a 128-byte granularity. \footnote{Although the PTX documentation specifies LDSM/STSM as operating on $8 \times 8$ matrices of 16-bit elements \cite[\S 9.7.13.4.15-16]{ptx}, we utilize these instructions for FP8 by packing pairs of 8-bit elements to fit this format. However, since the transpose variants of LDSM/STSM cannot unpack the 8-bit pairs, some register-level data shuffling must occur between LDSM and STSM to accomplish a tile-wise transpose; we omit those specifics here.} These instructions efficiently use registers, facilitating execution within the producer warpgroup, and inherently support layout transposition during memory copy. Additionally, beyond the initial iteration, the transpose for the next $\mathbf{V}$ tile can be scheduled to run concurrently—i.e., in the background—while the two WGMMAs that operate on the previous $\mathbf{V}$ and the present $\mathbf{K}$ tile are active.

Second, our analysis reveals that, in contrast to FP16, the organization of the FP32 accumulator associated with an FP8 WGMMA operation diverges from the memory arrangement expected for operand A when it is stored in registers. Illustrations of selected portions of both these register layouts are provided in \cref{fig:rmem_accum} and \cref{fig:rmem_operand}, with each thread possessing entries in the sequence shown. To align the accumulator output from the initial WGMMA with the requirements of a subsequent WGMMA—and to maintain compatibility with the structure of the $\mathbf{V}$ tile created by the in-kernel transpose—we employ byte permutation instructions. Specifically, as visualized in \cref{fig:rmem_accum}, the entries are reordered as
$$\{ \verb|d0 d1 d4 d5 d2 d3 d6 d7| \},$$
and this pattern is repeatedly applied for every 8 bytes segment. When considering the logical organization of the $\mathbf{P}$ tile, this effectively rearranges its columns (for instance, columns $0189$ now constitute the leading four columns). For the next WGMMA stage to generate the intended output tile, the in-kernel transpose can accordingly emit the $\mathbf{V}$ tile with a corresponding permutation of its rows.\footnote{This extra flexibility resulting from performing the in-kernel transpose obviates the necessity for shuffle instructions that would otherwise reassign register values between threads—a method we had previously discussed in~\citep{colfax_fp8_flashattention_2024}.}

\textbf{Accuracy: block quantization and incoherent processing.} The FP8 (e4m3) numerical representation allocates only 3 bits to the mantissa and 4 bits to the exponent, resulting in greater numerical imprecision in comparison to FP16 or BF16 formats. In addition, large-scale models are known to contain outlier values~\citep{dettmers2208llm, sun2024massive} whose magnitudes considerably exceed those of the majority of model parameters, creating challenges for the quantization process. Conventional quantization methods employ per-tensor scaling~\citep{micikevicius2022fp8}, which involves assigning a single scale factor to each tensor (such as one for $\mathbf{Q}$, one for $\mathbf{K}$, and another for $\mathbf{V}$).
To address the increased numerical errors in FP8 attention mechanisms, we apply two strategies:

\begin{enumerate}

\item \textbf{Block quantization}: Here, a single scalar is maintained for each block, meaning that for each of $\mathbf{Q}$, $\mathbf{K}$, $\mathbf{V}$, the tensor is partitioned into blocks of dimension $B_r \times d$ or $B_c \times d$, and quantization is applied to each block independently. This process can be combined with any operation occurring just prior to attention, such as rotary embedding, without introducing extra latency since rotary embedding is predominantly limited by memory bandwidth. Because the \textsc{FlashAttention-3} algorithm inherently processes data in block format, it is straightforward to rescale each block of $\mathbf{S}$ in line with the block quantization, all while incurring no additional computational overhead.
\item \textbf{Incoherent processing}: To mitigate the effects of outliers, $\mathbf{Q}$ and $\mathbf{K}$ are pre-multiplied by a random orthogonal matrix $\mathbf{M}$ prior to their quantization into FP8 format. Given the orthogonality of $\mathbf{M}$, where $\mathbf{M} \mathbf{M}^\top = I$, it follows that $(\mathbf{Q} \mathbf{M}) (\mathbf{K} \mathbf{M})^\top = \mathbf{Q} \mathbf{K}^\top$; thus, this transformation preserves the attention mechanism’s output. This approach functions by dispersing outliers, as each element of $\mathbf{Q} \mathbf{M}$ or $\mathbf{K} \mathbf{M}$ represents a randomly weighted sum of elements from the original $\mathbf{Q}$ or $\mathbf{K}$, which lessens quantization artifacts. Practically, we adopt the construction of $\mathbf{M}$ from \citet{chee2024quip} and~\citet{tseng2024quip}, wherein $\mathbf{M}$ consists of a sequence of random diagonal matrices with $\pm 1$ entries and a Hadamard matrix. This allows multiplication to be performed with $O(d \log d)$ complexity as opposed to $O(d^2)$, and such multiplication can also be merged with rotary embedding without incurring any additional computation time.
\end{enumerate}

Through experimental evaluation, we confirm in \cref{sec:numerical_error} that these two methods yield up to a 2.6$\times$ reduction in numerical error.

\section{Empirical Validation}
\label{sec:exp}

To develop \textsc{FlashAttention-3} and assess its performance and precision, we leverage the WGMMA and TMA abstractions provided by CUTLASS~\citep{Thakkar_CUTLASS_2023}.

\begin{itemize}

\item \textbf{Benchmarking attention.} The runtime performance of \textsc{FlashAttention-3} is evaluated across varying sequence lengths and compared against both the standard PyTorch implementation, \textsc{FlashAttention-2}, its Triton variant leveraging H100-specific instructions, and the cuDNN-based vendor-optimized \textsc{FlashAttention-2} for H100 GPUs. Our results show that \textsc{FlashAttention-3} delivers speedups of up to 2.0$\times$ over \textsc{FlashAttention-2} and up to 1.5$\times$ over its Triton implementation. In terms of throughput, \textsc{FlashAttention-3} achieves up to 740 TFLOPs/s, corresponding to 75\% of the peak theoretical TFLOPs/s attainable on H100 GPUs.
\item \textbf{Ablation study.} The acceleration observed in \textsc{FlashAttention-3} is attributable to the integration of warp-specialization and GEMM-softmax pipelining, as demonstrated by our ablation analyses of these algorithmic enhancements.
\item \textbf{Accuracy of FP8 attention.} Through empirical validation, we demonstrate that implementing block quantization and incoherent processing reduces numerical errors in FP8 computations for \textsc{FlashAttention-3} by a factor of 2.6$\times$.
\end{itemize}

\subsection{Benchmarking Attention}
\label{subsec:benchmark_attn}

We evaluate the execution times of various attention mechanisms on an H100 80GB SXM5 GPU across a range of configurations (with and without causal masking, and head dimensions of 64 or 128), using FP16 input data. The experimental findings, presented in~\cref{fig:benchmark_attn_fwd} and~\cref{fig:benchmark_attn_bwd}, demonstrate that \textsc{FlashAttention-3} achieves a 1.5-2.0$\times$ speedup over \textsc{FlashAttention-2} during the forward pass, and is 1.5-1.75$\times$ faster during the backward pass. Relative to the conventional attention implementation, \textsc{FlashAttention-3} delivers up to 3-16$\times$ improvement in speed. Notably, for sequence lengths of 1k and greater, \textsc{FlashAttention-3} outperforms even a proprietary GPU library (cuDNN -- closed source) that is specifically tuned for H100 hardware.

\paragraph{Benchmark settings:} We vary the sequence length as 512, 1k, ..., 16k, and set batch size so that the total number of tokens is 16k. We set the hidden dimension to 2048, and head dimension to be either 64, 128, or 256 (i.e., 32 heads, 16 heads, or 8 heads). To calculate the FLOPs of the forward pass, we use:

\begin{equation*}
  4 \cdot \text{seqlen}^2 \cdot \text{head dimension} \cdot \text{number of heads}.
\end{equation*}

Under causal masking, this quantity is halved since roughly 50\% of the elements are actually computed. For the backward pass, the number of FLOPs is estimated by multiplying the forward pass FLOPs by 2.5, reflecting that the backward computation involves 5 matrix multiplications compared to the 2 in the forward computation, as a result of recomputation.

Additionally, we evaluate the runtime performance of FP8 during the forward pass under comparable conditions. The outcomes for headdim 256 are depicted in ~\cref{fig:benchmark_attn_fp8}, while comprehensive findings are provided in \cref{sec:benchmark_attn_fp8_full}.

\subsection{Ablation Study: 2-Stage Pipelining Experiments}
\label{sec:2-stage-pipelining-experiments}

We conduct an ablation study on both the 2-stage WGMMA-softmax pipelining and the use of warp-specialization within non-causal FP16 \textsc{FlashAttention-3}, maintaining the following parameter settings: $\{\text{batch}, \text{seqlen}, \text{nheads}, \text{hdim}\} = \{ 4, 8448, 16, 128\}$. The outcomes, as presented in~\cref{table:ablation_pipelining}, demonstrate that incorporating our algorithmic enhancements—specifically, asynchrony introduced by warp-specialization and the overlapping execution of GEMM and softmax—yields substantial performance gains, improving throughput from 570 to 661 TFLOPs.

\subsection{Numerical Error Validation}
\label{sec:numerical_error}

Given the growing attention to the numerical error associated with \textsc{FlashAttention}~\citep{golden2024flash}, we conduct a comparison among \textsc{FlashAttention-2}, \textsc{FlashAttention-3}, and a conventional attention mechanism by benchmarking them against a reference implementation utilizing FP64 precision. To mimic the occurrence of outlier features and activations that are present in LLMs~\citep{dettmers2208llm, sun2024massive}, we sample the elements of $\mathbf{Q}, \mathbf{K}, \mathbf{V}$ from the following probability distribution:

\begin{equation*}
  \mathcal{N}(0, 1) + \mathcal{N}(0, 100) \cdot \mathrm{Bernoulli}(0.001).
\end{equation*}

In this setup, each matrix element follows a normal distribution with a mean of zero and a standard deviation of 1. However, for 0.1\% of the elements, we introduce an additional independent component sampled from a normal distribution with standard deviation 10. The root mean squared error (RMSE) outcomes are reported in~\cref{table:numerical_error}. Under FP16 precision, both \textsc{FlashAttention-2} and \textsc{FlashAttention-3} yield RMSE values that are 1.7$\times$ lower than those of the conventional method, owing to the preservation of intermediate softmax computations in FP32. In contrast, the FP8 baseline utilizes per-tensor scaling, accumulates matmul operations in FP32, and retains softmax intermediates in FP16. Owing to its use of block quantization and incoherent processing, \textsc{FlashAttention-3} executed in FP8 demonstrates a 2.6$\times$ improvement in accuracy relative to this baseline.

\section{Dicussion, Limitations, Conclusion}
\label{sec:discussion}

Through the development of \textsc{FlashAttention-3}, we illustrate how leveraging innovative programming strategies and hardware advancements—specifically asynchrony and low-precision operations—can substantially improve both the performance and fidelity of attention mechanisms. Our approach yields a 1.5-2.0$\times$ acceleration in attention computations relative to \textsc{FlashAttention-2}, and achieves a 2.6$\times$ reduction in FP8-related numerical errors when compared to conventional per-tensor quantization methods. There are several aspects that remain open for future enhancement, including optimizing performance for LLM inference scenarios, incorporating a persistent kernel architecture within the FP8 kernel,\footnote{In our benchmarks, the FP16 version of \textsc{FlashAttention-3} utilizes a persistent kernel along with a load-balancing technique, whereas the FP8 variant does not, partially accounting for the comparatively lower performance of FP8 \textsc{FlashAttention-3} at shorter sequence lengths and when using causal masking, relative to FP8 cuDNN kernels.} and further investigation into the implications of employing low-precision attention during large-scale model training. Although our evaluation is centered around Hopper GPUs, we anticipate that these methods will generalize to a variety of hardware accelerators. Ultimately, we envision that advances in the efficiency and accuracy of attention primitives such as those achieved here will pave the way for expanded capabilities in tasks involving extended context.

\end{document}